\subsection{ORTHOGONALITY}

DEFINITION 3. Let E be an A-module and \(\mathbf{E}^*\) its dual; an element \(x \in \mathbf{E}\) and an element \(x^* \in \mathbf{E}^*\) are called orthogonal if \(\langle x, x^* \rangle = 0\) .

A subset M of E and a subset \(M'\) of \(E^{*}\) are called orthogonal sets if, for all \(x \in M\) , \(x^{*} \in M'\) , x and \(x^{*}\) are orthogonal. In particular, \(x^{*} \in E^{*}\) (resp. \(x \in E\) ) is called orthogonal to M (resp. \(M'\) ) if it is orthogonal to every element of M (resp. \(M'\) ). If \(x^{*}\) and \(y^{*}\) are orthogonal to M, so is \(x^{*} + y^{*}\) and \(x^{*}\alpha\) for all \(\alpha \in A\) by virtue of (10) and (12) (no. 3), which justifies the following definition:

DEFINITION 4. Given a subset M of E (resp. a subset M' of E*), the set of \(x^* \in \mathbf{E}^*\) (resp. the set of \(x \in \mathbf{E}\) ) which are orthogonal to M (resp. M') is called the submodule totally orthogonal to M (resp. M') (or simply the submodule orthogonal to M (resp. M') if no confusion can arise).

By definition of a linear form, the submodule of \(E^{*}\) orthogonal to E is reduced to 0; the submodule of \(E^{*}\) orthogonal to \(\{0\}\) is identical with \(E^{*}\) .

PROPOSITION 6. Let M, N be two subsets of E such that M ⊂ N; if M' and N' are the submodules of E* orthogonal to M and N respectively, then N' ⊂ M'.

PROPOSITION 7. Let \((\mathbf{M}_t)\) be a family of subsets of E; the submodule orthogonal to the union of the \(\mathbf{M}_t\) is the intersection of the submodules \(\mathbf{M}_t'\) which are respectively orthogonal to the \(\mathbf{M}_t\) ; this submodule is also the submodule orthogonal to the submodule of E generated by the union of the \(\mathbf{M}_t\) .

These results are immediate consequences of the definitions.

There is an analogous proposition (which we shall leave to the reader to state) for submodules of E orthogonal to subsets of \(E^{*}\) .

If M is a submodule of E, \(M'\) the submodule of \(E^{*}\) orthogonal to M and \(M''\) the submodule of E orthogonal to \(M'\) , then \(M \subset M''\) but it may be that \(M \neq M''\) (Exercise 9). Note however that if \(M''\) is the orthogonal of \(M''\) in \(E^{*}\) , then \(M'' = M'\) ; for \(M' \subset M''\) and on the other hand the relation \(M \subset M''\) implies \(M'' \subset M'\) .

